% Auto-generated by paper/build_survey.py -- do not edit by hand.
{\footnotesize
\setlength{\tabcolsep}{3pt}
\renewcommand{\arraystretch}{1.15}
\begin{longtable}{@{}p{2.1cm} >{\raggedright\arraybackslash}p{1.9cm} >{\raggedright\arraybackslash}p{1.5cm} >{\raggedright\arraybackslash}p{1.3cm} >{\raggedright\arraybackslash}p{1.5cm} >{\raggedright\arraybackslash}p{2.4cm} >{\raggedright\arraybackslash}p{3.0cm}@{}}
\caption{Processing choices of the 103 surveyed \dvv\ studies (the literature survey underpinning this paper). Every study is cited; \texttt{n/r} = not reported in the source. Frequency band, coda window, estimator, reference/stack scheme and uncertainty treatment are the choices Section~\ref{sec:results} shows to control the result.}\label{tab:survey}\\
\toprule
\textbf{Study} & \textbf{Site / region} & \textbf{Freq (Hz)} & \textbf{Coda (s)} & \textbf{Method} & \textbf{Reference / stack} & \textbf{Uncertainty treatment} \\
\midrule
\endfirsthead
\multicolumn{7}{@{}l}{\footnotesize\itshape Table~\ref{tab:survey} continued}\\
\toprule
\textbf{Study} & \textbf{Site / region} & \textbf{Freq (Hz)} & \textbf{Coda (s)} & \textbf{Method} & \textbf{Reference / stack} & \textbf{Uncertainty treatment} \\
\midrule
\endhead
\midrule\multicolumn{7}{r@{}}{\footnotesize\itshape continued on next page}\\
\endfoot
\bottomrule
\endlastfoot
\citet{Brenguier2008} & Piton de la Fournaise, La Réunion & 0.1-0.9 & n/r (slope of dtau vs tau; short scanning window) & MWCS (doublet) & Daily stacks referenced to long-term stack & linear-slope uncertainty; exclude >0.04\% \\ \hline
\citet{Duputel2009} & Piton de la Fournaise, La Réunion & 0.1–1 & n/r & Other & Quasi-real-time stacks vs reference & n/r \\ \hline
\citet{Obermann2013} & Piton de la Fournaise, La Réunion & 0.1-1.0 & 10-45 & stretching & Daily stacks, referenced; coda at increasing lapse times & Weaver et al. coherence-based error formula \\ \hline
\citet{HotovecEllis2014} & Mount St. Helens, WA, USA & 1-10 (also 1-5, 5-10) & windows after S arrival; >=5 with CCC>0.65 & doublet (Snieder 2002) & Stack across several hundred repeating-earthquake families (multiplets) & \textasciitilde{}0.1\% from SD of velocity fit \\ \hline
\citet{HotovecEllis2015} & Mount St. Helens, WA, USA & 0.5-10 & 15 s window starting 2 s after first motion & stretching & Combines repeating-earthquake CWI with ambient-noise interferometry & reduced similarity (CC) in stretching window \\ \hline
\citet{Rivet2015} & Piton de la Fournaise, La Réunion & 0.25-2 & n/r & MWCS & Daily stacks vs reference; network-averaged dv/v & coherency + linear-regression error \\ \hline
\citet{DePlaen2016} & Piton de la Fournaise, La Réunion (method also re Kawah Ijen) & 0.1-1.0, 0.5-1.0, 1.0-2.0, 2.0-4.0 & 5-35 (both branches) & MWCS & n/r & coherency + linear-regression error QC \\ \hline
\citet{Donaldson2017} & Kīlauea summit, Hawaii, USA & 0.33-1.0 & 30 s window, min lag = interstation dist / 0.8 km/s & MWCS (doublet) & Daily NCFs stacked over 3-day moving window & weighted dt-t regression; coherence>0.65, err<0.1 s thresholds \\ \hline
\citet{Takano2017} & Izu-Oshima, Japan & 0.5-1, 1-2, 2-4 & -20 to +20 & MWCS & Cross-correlations 2012–2015, multiple bands & error bars from coherency \\ \hline
\citet{Lesage2018} & Volcán de Colima, Mexico & 0.125-2 & 10-80 & stretching & Daily cross-correlations; 2013 stack as reference & empirical noise from AVV fluctuations (\textasciitilde{}0.05\%) \\ \hline
\citet{DePlaen2019} & Mt. Etna, Italy (2013–2014) & 1.0-2.0 & 5-35 & MWCS & Daily autocorrelations, 2-day linear stack & cross-coherence + squared misfit; reject dt err>0.1 s or coh<0.6 \\ \hline
\citet{Donaldson2019} & Northern Volcanic Zone (Askja/Bárðarbunga), Iceland & 0.1-0.4, 0.4-1.0, 1-2, 2-4, 4-16 & per-band (suppl. Table S1) & stretching + MWCS & n/r & reject dv/v when stretched-ref CC<0.4 \\ \hline
\citet{Olivier2019} & Kīlauea, Hawaii, USA (2018 eruption) & 0.08-1.2 & 30 s windows; start = dist/700 m/s + 30 s & MWCS & n/r (passive image interferometry, daily) & CCF-coherence-based uncertainty \\ \hline
\citet{Yates2019} & Whakaari / White Island, New Zealand & 0.1-1.0 & 16 s moving windows within 20-80 s & MWCS & n/r & QC (err<0.1 s, coh>0.7); error-weighted mean \\ \hline
\citet{Feng2020} & Kīlauea, Hawaii, USA & 3-8 & [-2.8:-0.4] and [0.4:2.8] (=[-14Dt:2Dt],[2Dt:14Dt]) & stretching & Jan 2017–Jun 2018 NCFs & absolute error from stretching CC \\ \hline
\citet{HotovecEllis2022} & Kīlauea, Hawaii, USA (2018 collapse) & n/r & n/r & CWIRE (stretching + damped LSQ) & Coda Wave Interferometry with Repeating Earthquakes (CWIRE), per-cluster pairs & n/r \\ \hline
\citet{Kopfli2024} & Mount St. Helens, WA, USA & n/a (wavefield amplitude features, not dv/v) & n/a & n/a (no dv/v; companion is Makus 2024) & Long-term NCF stacks; many dv/v series on a common grid & n/a \\ \hline
\citet{Yates2024} & Mt. Ruapehu, New Zealand (2005–2009) & 0.25-2.5 (wavelet 0.1-8.0) & min lag 5 s, length 20 cycles (e.g. 5-25 at 1 Hz) & CWT (cross-wavelet) & n/r & slope SD from covariance; coherence weighting \\ \hline
\citet{Yukutake2025} & Izu-Oshima, Japan (2003–2020) & 0.1-0.9, 0.5-2.0, 1.0-4.0 & 20-40 & MWCS & Long-term interferometry stacks, 2003–2020 & 1-2 sigma errors as inversion weights \\ \hline
\citet{Schaff2004} & 1989 Mw 6.9 Loma Prieta aftershock zone, California (also Parkfield repeaters) & n/r (unfiltered) & 1.4 s moving windows through coda & doublet & Repeating-earthquake doublets; delays vs pre-mainshock repeaters & formal slope SE; CC weighting \\ \hline
\citet{Pacheco2005} & Theory / acoustic diffusion model & n/r & n/r & Other & n/r & n/r (analytical kernel) \\ \hline
\citet{Rubinstein2005} & 2004 M6 Parkfield earthquake, San Andreas fault, California & n/r & n/a (direct-S delays, not coda) & doublet (moving-window CC) & Repeating-earthquake travel-time delays before vs after mainshock & CC>0.8 and SNR>4:1 QC \\ \hline
\citet{Wegler2007} & Mid-Niigata (Chuetsu) earthquake source region, Japan; F-Net station KZK (\textasciitilde{}24 km) & 2-100 (2 Hz high-pass) & 5-14 & stretching & Daily autocorrelations vs \textasciitilde{}two-week pre-event reference; grid search (10000 trials, dv/v… & day-to-day fluctuation \textasciitilde{}0.1\%; CC>0.5 \\ \hline
\citet{Brenguier2008b} & Parkfield, San Andreas fault, California, USA & 0.1–0.9 & n/r & Stretching & 1550-day reference vs overlapping 5-day moving windows; 30-day time resolution & Averaged over 78 receiver pairs plus 5-day segments \\ \hline
\citet{Hadziioannou2009} & Laboratory scattering gel (ultrasonic) & 2.5 MHz ultrasonic (lab, not seismic band) & 12.5-50 us (lab) & stretching + doublet (compared) & Noise-correlation stacking & analytic std of stretching-CC fluctuations \\ \hline
\citet{Sawazaki2009} & Western-Tottori region, Japan; KiK-net borehole (surface + 100 m); 2000 Mw 6.6 mainshock & n/r & n/r & Other & Coda deconvolution of surface vs 100-m borehole records & n/r \\ \hline
\citet{Wegler2009} & 2004 Mw 6.6 Mid-Niigata earthquake source region, Japan; 6 Hi-net/F-Net stations <25 km & 2-8 and 0.1-0.5 & 5-14 (2-8 Hz); 10-100 (0.1-0.5 Hz) & stretching & Daily auto-/cross-correlations inverted daily for 2 months pre/post; pre-event reference & SD over Green's-tensor components; R>0.5 \\ \hline
\citet{Chen2010} & Sichuan, China — 2008 Mw 7.9 Wenchuan earthquake fault zone & n/r & \textasciitilde{}25 s after ballistic, up to +-200 s lapse & MWCS & Noise cross-correlation stacking; sub-array comparison (156 broadband stations) & n/r \\ \hline
\citet{Nakata2011} & NE Japan (Honshu), KiK-net; 2011 Tohoku-Oki earthquake & 1-13 & n/a (deconvolution interferometry, direct S-wave) & n/a (not coda/noise dv/v) & >300 earthquakes (Jan 1–May 26 2011); borehole-to-surface deconvolution & SD of travel times over events \\ \hline
\citet{Rivet2011} & Guerrero, Mexico subduction zone (2006 M7.5 slow slip event) & 0.037-0.27 (periods 3.7-27 s) & n/r (coda used, no length) & doublet (MWCS) & n/r & previous-epoch referencing; no explicit error bars \\ \hline
\citet{Hobiger2012} & 2008 Iwate-Miyagi Nairiku earthquake (Mw 6.9), NE Japan & 0.125-0.25, 0.25-0.5, 0.5-1.0 & ten periods, starting 7.5 periods after direct wave & stretching & Daily cross-correlations vs long-term reference Green's function & CC weighting; 18-component weighted average \\ \hline
\citet{Minato2012} & Southern Tohoku (Fukushima/Ibaraki), Japan; 58 Hi-net stations; 2011 Tohoku-Oki & 2-5 & 2-10 lag & stretching & n/r & SD from pre/post-seismic values \\ \hline
\citet{Obermann2013} & Numerical 2-D elastic media (lunar-data illustration) & 20 Hz central, \textasciitilde{}12 Hz bw (numerical) & 1.5 s windows, centered \textasciitilde{}1.8-6.6 s & stretching & Averaged over 10 random medium realizations per configuration & std over ten random-media realizations \\ \hline
\citet{Brenguier2014} & Volcanic/geothermal regions of Honshu, Japan; 2011 Mw 9.0 Tohoku-Oki & 0.1–0.9 & n/r & MWCS & Daily cross-correlations; linear inversion/regularization for continuous time series & Stress sensitivity \textasciitilde{}0.001 MPa\textasciicircum{}-1 inferred; detailed dv/v uncertainty n/r \\ \hline
\citet{Liu2014} & Epicentral region of 2008 Mw 7.9 Wenchuan earthquake, China & 0.125-1 (1-8 s; subbands 1-2,2-4,4-8 s) & 30-130 & MWCS & Noise cross-correlation stacking (Aug 2004 – Sep 2011) & 2-sigma SD across nine components \\ \hline
\citet{Taira2015} & South Napa, California (2014 Mw 6.0 South Napa earthquake) & 0.1-0.9 & -60 to -20 and +20 to +60 & MWCS & MSNoise cross-correlations; decimated to 20 Hz & Jackknife 95\% CI; 2-sigma SD \\ \hline
\citet{Gassenmeier2016} & Station PATCX, Atacama Desert, Northern Chile (IPOC) & 1-3, 4-6, 7-10 (primary 4-6) & 10-15 (for 4-6 Hz) & stretching (autocorrelation CWI) & Daily autocorrelations; two-step iterative reference & Gaussian fit to correlation values (CIs) \\ \hline
\citet{Hillers2019} & San Jacinto fault zone, California (after 2010 M7.2 El Mayor-Cucapah, M5.4 Collins Valley) & 0.2-2 (sub-bands 0.2-0.4,0.4-0.8,0.8-1.6,0.3-1.5) & 20-50 (also 20-40,30-50,40-60) & stretching and MWCS (doublet) & Average-waveform reference for stretching; MWCS also applied & P-value QC; cc>=0.85 inversion \\ \hline
\citet{Wang2019} & NE Honshu, Japan (2011 Mw 9.0 Tohoku-Oki earthquake) & period bands 8-30 s and 15-50 s & -400 to 400 (moving 30 \& 50 s) & doublet & Monthly velocity-change estimates & SD error bars from 2009-2010 \\ \hline
\citet{Mao2020} & Salton Sea Geothermal Field, California (2009–2011) & 0.5-2.0, 0.75-3.0, 1.1-2.2, 1.5-6.0, 2.0-8.0 & n/r (real data); synthetic 18.65-35 & wavelet cross-spectrum (WCS) & Noise cross-correlation stacking & wavelet cross-spectrum amplitude weighting; coherence threshold \\ \hline
\citet{Poli2020} & L'Aquila region, central Italy (2009 Mw 6.3 L'Aquila earthquake) & 0.5-1 & start 10/20/30 s, length 20/40/60 s & stretching (doublet tested) & n/r & CC-weighting; cc>0.9 retained \\ \hline
\citet{Boschelli2021} & Ridgecrest fault zone, California (2019 Mw 7.1 Ridgecrest earthquake) & >1 (1 Hz high-pass) & n/r (variable lapse window) & stretching & Daily autocorrelation functions vs mean waveform & SD across non-overlapping windows \\ \hline
\citet{Lu2021} & Ridgecrest, California (2019 Mw 7.1 Ridgecrest earthquake) & 8.0–12.0 & 3 (moving windows, 1.5 s step) & Other & 10-min stacks; adaptive Gaussian smoothing (20 min growing to 24 hr over 3 days) & Standard deviations tracked; stabilized via smoothing \\ \hline
\citet{Sheng2022} & San Jacinto fault, Anza seismic gap, Southern California & >1 & n/r & wavelet (WTS) & Weekly stacked correlations with \textasciitilde{}2-month smoothing & dt uncertainty as inverse weight \\ \hline
\citet{Mainsant2012} & Pont-Bourquin landslide, Swiss Alps (Switzerland) & 10-12 (analyzed 4-25) & 0.2-2 (both sides) & stretching & Daily noise correlograms (cross-correlation functions) & Weaver et al. (2011) CC-based error (\textasciitilde{}1-2\%) \\ \hline
\citet{Voisin2016} & Avignonet / Mas d'Avignonet landslide, Trièves (French Alps), France & \textasciitilde{}6-8 (annual dv/v pattern best resolved) & n/r & Stretching & Daily/seasonal stacks vs reference & Correlation-coefficient weighting; comparison with piezometer \\ \hline
\citet{Harba2017} & Just-Tęgoborze landslide, Carpathian flysch, southern Poland & High-frequency seismic noise (\textasciitilde{}5-20+, dispersion-curve range) & n/r (interferometric Green's-function retrieval) & Other & Cross-correlation / interferometric stacks; dispersion inversion (neighbourhood algorithm) & Dispersion-inversion misfit; comparison with MASW \\ \hline
\citet{Bertello2018} & Montaguto earthflow, southern Italy & n/a (active-source ReMi/MASW dispersion, not coda dv/v) & n/a & n/a & Time-lapse interferometry + time-lapse active MASW & n/a \\ \hline
\citet{Bievre2018} & Pont-Bourquin landslide, Swiss Alps (4.5-yr record) & \textasciitilde{}5-15 (surface-wave band; reported sensitivity to shallow layer <=2 m) & n/r & Stretching & Daily correlations, multi-year reference & Correlation coefficient; seasonal-cycle modeling to separate from precursors \\ \hline
\citet{Colombero2018} & Madonna del Sasso cliff, NW Italian Alps (Orta Lake), Italy & 2-20 (analysis); 2-4 strongest annual signal & 0.5-2 (and symmetric -2 to -0.5) & Stretching & Daily correlations, multi-year reference (2013-2016) & Correlation coefficient (\textasciitilde{}0.9 at low freq) as reliability metric \\ \hline
\citet{Bontemps2020} & Maca slow-moving Andean landslide, southern Peru (3-yr dataset) & n/r & n/r & n/r & Daily stacks vs reference & n/r \\ \hline
\citet{Fiolleau2020} & Harmalière landslide, French Western Alps (collapsed Nov 2016) & 1-12 (monitoring); dv/v in 2-4 and 8-10 bands; block resonance \textasciitilde{}9-16 & 0.05-1.5 & Stretching & Hourly records, daily averaging vs reference & Correlation-coefficient degradation tracked across bands; multi-parameter precursor timing \\ \hline
\citet{LeBreton2021} & Global review (9 landslides monitored to date) & 1.3-20 (commonly 10-14; 0.1-1 for large/deep) & end of direct waves to 5-10 periods (site-dependent) & stretching and doublet (MWCS) & Daily stacks vs long-term reference & Weaver et al. (2011) rms + SD over pairs/windows; discard CC<0.6 \\ \hline
\citet{Fan2023} & Rock slope, southwest China & High-frequency (\textasciitilde{}5-20) & Sub-second to a few s & Stretching & Short-window stacks for rapid detection & Decorrelation/CC checks; rapid-detection confidence discussion \\ \hline
\citet{Liu2025} & Xishan Village landslide, Sichuan Province, China & 5-25 (raw 1-30); 15 Hz top \textasciitilde{}15 m, 5 Hz 15-40 m & 0.4-2.0 & Stretching & 20-min interval CCFs vs reference (SNR>4 wave-packet selection) & 95\% confidence interval (\textasciitilde{}0.05\% uncertainty at 5 Hz) \\ \hline
\citet{Whiteley2026} & Hollin Hill Landslide Observatory, North Yorkshire, UK (2-yr) & 6-10 (plus two unspecified bands) & n/r & MWCS & Time-lapse interferometry + MASW between recording periods & none (data report) \\ \hline
\citet{deWit2026} & Open-pit mine slope, Australia & n/r (high-frequency near-surface) & n/r & Stretching & Time-lapse stacks vs reference & Decorrelation/CC; comparison with radar surface deformation \\ \hline
\citet{SensSchonfelder2006} & Merapi Volcano, Indonesia & 0.5 Hz high-pass & 2-8 & stretching & Daily autocorr vs yearly reference & SD over non-overlapping windows \\ \hline
\citet{Tsai2011} & Theoretical model, applied to southern California & n/a (theory of thermoelastic/hydrologic velocity change) & n/a & n/a (analytic model, no measurement) & n/r & n/a \\ \hline
\citet{Hillers2014} & TCDP borehole array, Chelungpu fault, Taiwan & n/r & n/r & n/r & n/r & n/r \\ \hline
\citet{Lecocq2017} & Gräfenberg Array (GRA1-4), SE Germany (karst limestone aquifer) & 0.1-0.8 & 20-100 (and -20 to -100) & MWCS (Brenguier all-pairs) & Daily CCFs, 31-day rolling window & LSQ slope error as inverse-variance weight; bootstrap \\ \hline
\citet{Nimiya2017} & Kyushu Island, Japan (2016 Kumamoto region) & 0.1-0.9 & 100 & stretching (+MWCS check) & Daily CCFs from 30-min segments; 1-yr reference, 30-day moving stack & SD over six 50-s subwindows \\ \hline
\citet{Wang2017} & Japan (nationwide; strongest in Kyushu/volcanic zones) & 0.1-0.25 and 0.5-2.0 & 30 s (0.1-0.25 Hz), 10 s (0.5-2 Hz) & MWCS & Daily CCFs; \textasciitilde{}1-yr reference, moving-window current stack & none (QC only: dt<=0.2 s, coh>=0.5) \\ \hline
\citet{Clements2018} & San Gabriel Valley, California, USA & 0.1-0.25 and 0.5-2.0 & 30 s (0.1-0.25 Hz), 10 s (0.5-2 Hz) & MWCS & Daily stacks vs reference & none (QC only: dt<=0.2 s, coh>=0.5) \\ \hline
\citet{Kim2019} & Gulf Coast Aquifer, southern Texas (IU.HKT near Houston), USA & 0.01-8 & n/r & stretching & Multi-year (\textasciitilde{}20 yr) record; yearly-scale comparisons & none (grid-search CC max) \\ \hline
\citet{Andajani2020} & Chugoku \& Shikoku, SW Japan (Hi-net) & 0.1-0.9 & 100 & stretching & n/r & stretching CC as quality indicator (no error bars) \\ \hline
\citet{GaubertBastide2022} & Crépieux-Charmy water-production field, Lyon, France & 2-5 & n/a (ballistic Love-wave window, not coda) & stretching & Hourly correlations over 19 days, two fill/drain cycles & eps\_max +- sigma (stretching variance) \\ \hline
\citet{Illien2022} & Nepal Himalaya (2015 Mw7.8 Gorkha aftermath) & 4-8 (Chaku); 2-4 (Gumba) & 12 periods (\textasciitilde{}3 s) after 4-period skip & stretching & \textasciitilde{}3-yr continuous time series & none (multi-reference stack) \\ \hline
\citet{Mao2022} & Coastal Los Angeles basins, California, USA & 2-4 & 2-8 & stretching & Daily CCFs stacked over 20 days, 5-day step; pairs <50 km & cc\textasciicircum{}2-weighted channel mean (no explicit error bars) \\ \hline
\citet{Clements2023} & Statewide California, USA (1999-2021) & 2-4 & 2-8 & stretching & Daily stacks across \textasciitilde{}700 stations vs reference & cc\textasciicircum{}2-weighted channel mean (no explicit error bars) \\ \hline
\citet{Delouche2023} & Greece (aquifer monitoring sites) & 0.33-1 (1-3 s period) & 15-55 & stretching & n/r & Weaver et al. (2011) RMS error \\ \hline
\citet{Ermert2023} & Mexico City basin / Valley of Mexico & 0.5-1, 1-2, 2-4, 4-8 & 4-10x and 8-20x longest period (e.g. 8-16 \& 16-40 s) & stretching & Clustered (GMM) autocorrelation stacks & CC\_best>0.6 QC; no formal dv/v error \\ \hline
\citet{Fokker2023} & Groningen, The Netherlands & 1.3-1.6 & n/r (coda of cross-coherence) & stretching & n/r & standard error sigma/sqrt(n) \\ \hline
\citet{Zhang2023} & Central Oklahoma, USA & 0.1-1.0 (also 0.3-1.2, 0.5-1.5, 1.0-2.0) & n/r (dynamic windows, 3.0 \& 2.0 km/s) & MWCS & Continuous 2013-2022 record & SD (avg \textasciitilde{}0.009\%) \\ \hline
\citet{Mao2025} & Greater Los Angeles, California, USA & 0.1-0.3 & 300 s window from 1.3*t0 & stretching & Daily CCFs stacked, multi-day step; two-decade record & Weaver et al. (2011) error formula \\ \hline
\citet{Mordret2016} & Southwest Greenland ice sheet & 0.1-0.3 & 300 s window from 1.3*t0 & stretching & Daily cross-correlations over \textasciitilde{}2-year record & Weaver et al. (2011) error formula \\ \hline
\citet{James2019} & Interior Alaska (Poker Flat / Fairbanks area), USA & 3-30 (18 bands; 13 Hz inversion) & +-0.2 to 2.0 (center lag) & MWCS & Daily cross-correlations stacked; reference relative to Jan 2014 & 95\% CI from posterior \\ \hline
\citet{Guillemot2020} & Gugla rock glacier, Valais, Switzerland & n/a (resonance-freq tracking, not dv/v) & n/a & n/a (no dv/v computed) & Daily correlograms from hourly raw data & n/a \\ \hline
\citet{Guillemot2021} & Laurichard (French Alps) and Gugla (Swiss Alps) rock glaciers & n/a (resonance-freq tracking, not dv/v) & n/a & n/a (no dv/v computed) & PSD-based; finite-element modal modeling & n/a \\ \hline
\citet{Lindner2021} & Mt. Zugspitze, German/Austrian Alps & 2-8 (analysis 1-20) & -5 to 5 lag & wavelet cross-spectrum (Morlet) & 15-year continuous record; daily/seasonal stacks & 1 SD \\ \hline
\citet{Luo2023} & Greenland Ice Sheet (GLISN stations) & 0.1-1 (also 0.1-0.5, 1-2) & 20-70 & MWCS & Daily autocorrelations stacked & fitted +-sigma; 95\% CI \\ \hline
\citet{Gassenmeier2015} & Ketzin CO2 storage site (CO2SINK), Brandenburg, Germany & 1.5-3 & moving windows after 300 m/s phase & stretching & 1-hour segments cross-correlated, stacked to daily & Weaver et al. (2011); scaling-coeff CIs \\ \hline
\citet{Hillers2015} & Basel deep geothermal (EGS) reservoir, Switzerland & n/r & Coda of noise correlations & Stretching & Daily noise correlations around 2006 stimulation & Sensitivity-kernel imaging of velocity-change location \\ \hline
\citet{Obermann2015} & St. Gallen geothermal site, Switzerland & 0.1-1 & 20 s window centered \textasciitilde{}15 s in coda & stretching & Daily cross-correlations; reference stacks & none \\ \hline
\citet{Czarny2016} & Underground coal mine, Upper Silesia, Poland & 0.6-1.2 & n/r & Stretching & Continuous cross-correlations over \textasciitilde{}42 days & n/r \\ \hline
\citet{Olivier2017} & Active tailings storage facility (mine), South Africa & 0.25-1.0, 0.5-2.0, 0.75-3.0, 1.0-4.0, 1.5-6.0, 2.0-8.0 & 20 s (-40 to -20 and 20 to 40) & MWCS (MSNoise) & Daily cross-correlations across array & two-sigma SD; CC>0.85 rejection \\ \hline
\citet{Taira2018} & Salton Sea Geothermal Field, California, USA & 0.25-1.0, 0.5-2.0, 0.75-3.0, 1.0-4.0, 1.5-6.0, 2.0-8.0 & 20 s (-40 to -20 and 20 to 40) & MWCS (MSNoise) & Daily NCFs from 30-min segments; 5-day stacks & two-sigma SD; CC>0.85 rejection \\ \hline
\citet{Kristjansdottir2019} & Hellisheidi geothermal field, SW Iceland & n/r & n/r & Stretching & Daily correlations (MSNoise workflow) & Seasonal and noise-source effects discussed as caveats \\ \hline
\citet{Snieder2002} & theory + ultrasonic lab & n/a (dispersion/preprocessing paper, no dv/v) & n/a & n/a (not a velocity-change study) & n/r & n/a \\ \hline
\citet{Bensen2007} & USArray test data & 0.1-0.9 & 6 s windows overlapping 3 s & MWCS & daily cross-correlation + temporal stacking & WLS phase-fit error (corrected \textasciitilde{}6x) \\ \hline
\citet{Clarke2011} & synthetic + real noise tests & 0.1-0.9 & 6 s windows overlapping 3 s & MWCS & current stack vs reference, weighted linear regression dt vs t & WLS phase-fit error (corrected \textasciitilde{}6x) \\ \hline
\citet{Hadziioannou2011} & urban noise autocorrelations & 0.1-0.9 & n/r (whole coda t1-t2) & stretching (and doublet) & waveform clustering to enhance daily stacks & remnant coherence (CC) \\ \hline
\citet{Weaver2011} & theory & 0.05-0.5 & up to t0=30 & stretching & reference vs current correlation & max CC>0.9 (no formal error) \\ \hline
\citet{Zhan2013} & synthetic + field examples & 0.05-0.5 & up to t0=30 & stretching & reference vs current & max CC>0.9 (no formal error) \\ \hline
\citet{Lecocq2014} & Piton de la Fournaise (validation) & synthetic (Ricker 12.5; noise 0.5-30) & 1-10 lapse; sub-windows 0.1-1.0 s & windowed CC, stretching, DTW (compared) & daily CCF + reference stack & WLS lag-vs-lapse; CC>=0.6 \\ \hline
\citet{Mikesell2015} & numerical CWI examples & 0.33-1.0 (1-3 s) & (dist/1.8)+20 to 120 & stretching & n/a & Weaver et al. (2011) rms formula \\ \hline
\citet{Stehly2015} & Wenchuan Mw 7.9 region (test dataset) & 0.1-1.0 (real); 0.15-0.65 (synthetic) & [15,35] and [-35,-15] & stretching (+MWCS compare) & daily correlations curvelet-filtered before stacking & CC>0.7 rejection; \textasciitilde{}+-0.1\% significance \\ \hline
\citet{Daskalakis2016} & synthetic + real noise & n/r (numerical study) & 1 s overlapping windows & stretching & normalized cross-correlation functions & SD over 20 receivers \\ \hline
\citet{Obermann2016} & 3-D wavefield simulations & 0.15-0.90 & -60 to +60 (doublet over coda) & doublet + Bayesian least-squares inversion (Brenguier 2014) & n/a & Bayesian least-squares inversion; error bars from inversion \\ \hline
\citet{Obermann2019} & review & scale-dependent & lapse-time/depth dependent & Other & reference and moving-stack strategies reviewed & synthesizes error sources: noise-source variability, processing choices, sensitivity ker… \\ \hline
\citet{Jiang2020} & benchmark vs MSNoise & user-defined & user-defined & Other & linear/PWS/robust stacking; reference + moving stacks & offers multiple dv/v estimators so users can cross-check; parallel/HDF5 reproducible pip… \\ \hline
\citet{Wang2020} & review & n/r (review) & n/r (review) & stretching, MWCS, DTW, WCS (review) & reference and stacking strategies reviewed & Bayesian least-squares / MCMC (review) \\ \hline
\citet{Yuan2021} & 2-D heterogeneous half-space simulations & varies (e.g. 0.5-2.2) & 45-75 (example) & WCC, TS, DTW, MWCS, WCS, WTS, WTDTW (7 compared) & n/a & residuals vs truth over realizations \\ \hline
\end{longtable}
}
